\chapter{Muscle Rigidity}

\section*{Definition}
Muscle rigidity is a hypertonic state characterized by increased resistance to passive movement that is velocity-independent, distinguishing it from spasticity. It affects both flexors and extensors uniformly and may be described as cogwheel (with superimposed tremor) or lead-pipe (smooth, continuous).

\section*{What Happens in Your Body}
Rigidity arises from dysfunction of the basal ganglia motor circuit, specifically from reduced dopamine in the substantia nigra pars compacta leading to increased activity of the indirect pathway. This results in excessive inhibition of the thalamus and over-excitation of the corticospinal tract. Loss of cortico-striatal inhibition produces sustained contraction of agonist and antagonist muscles simultaneously. Cogwheel rigidity reflects an underlying 4--6 Hz resting tremor superimposed on hypertonia.

\section*{Causes}
\begin{itemize}
  \item \textbf{Parkinson disease:} Most common cause; asymmetric, cogwheel rigidity with bradykinesia and resting tremor
  \item \textbf{Atypical parkinsonism:} Multiple system atrophy (MSA), progressive supranuclear palsy (PSP), corticobasal degeneration (CBD), dementia with Lewy bodies (DLB)
  \item \textbf{Drug-induced:} Neuroleptics (antipsychotics), metoclopramide, lithium, valproate, MPTP
  \item \textbf{Toxin-induced:} Carbon monoxide poisoning, manganese toxicity, cyanide, methanol
  \item \textbf{Metabolic:} Wilson disease, hypoparathyroidism, nonketotic hyperglycemia
  \item \textbf{Infectious:} Viral encephalitis, HIV-associated parkinsonism
  \item \textbf{Structural:} Basal ganglia stroke, tumor, trauma, hydrocephalus
\end{itemize}

\section*{When to See a Doctor}
See your doctor if:
\begin{itemize}
  \item New-onset rigidity with bradykinesia, resting tremor, or postural instability
  \item Rigidity with altered mental status, fever, or autonomic instability (rule out neuroleptic malignant syndrome)
  \item Rapid progression over weeks to months
  \item Onset after initiating antipsychotic medication
\end{itemize}

\section*{Related Symptoms}
\begin{itemize}
  \item Bradykinesia (slowness of movement)
  \item Resting tremor (pill-rolling, 4--6 Hz)
  \item Postural instability and festinating gait
  \item Masked facies (hypomimia)
  \item Micrographia, hypophonia, and drooling
\end{itemize}
\textbf{Important:} Acute rigidity with hyperthermia and altered mental status after neuroleptic exposure is neuroleptic malignant syndrome and requires emergency treatment.

\section*{Self-Care / Home Management}
\begin{itemize}
  \item Daily stretching exercises focusing on neck, trunk, and extremities
  \item Fall prevention: remove trip hazards, install grab bars, use walking aids
  \item Large-print materials and adaptive devices for micrographia
  \item Adequate hydration and high-fiber diet for constipation (common in parkinsonism)
  \item Over-the-counter laxatives as needed for bowel regularity
\end{itemize}

\section*{How Your Doctor Will Evaluate You}
\begin{itemize}
  \item Detailed neurological examination including UPDRS (Unified Parkinson Disease Rating Scale)
  \item Response to levodopa challenge (significant improvement suggests Parkinson disease)
  \item MRI brain to exclude structural lesions, vascular parkinsonism, or atypical syndromes
  \item Dopamine transporter (DaT) SPECT scan to differentiate Parkinson disease from tremor disorders
  \item Laboratory studies: TSH, copper, ceruloplasmin, HIV, serum protein electrophoresis
\end{itemize}

\section*{Treatment Options}
\begin{itemize}
  \item Levodopa/carbidopa as the most effective symptomatic treatment
  \item Dopamine agonists (pramipexole, ropinirole) for early disease or as adjunct
  \item MAO-B inhibitors (selegiline, rasagiline) for mild symptoms
  \item Anticholinergics (benztropine) for rigidity in younger people
  \item Deep brain stimulation of the subthalamic nucleus or globus pallidus for advanced disease
  \item Withdrawal of causative medications in drug-induced rigidity
  \item Dantrolene and supportive care for neuroleptic malignant syndrome
\end{itemize}

\section*{References}
\begin{thebibliography}{99}
\bibitem{muscle_rigidity:ref1} Jankovic J. "Parkinson's disease: clinical features and diagnosis." \textit{Journal of Neurology, Neurosurgery \& Psychiatry}, 2008;79(4):368--376.
\bibitem{muscle_rigidity:ref2} Postuma RB, Berg D, Stern M, et al. "MDS clinical diagnostic criteria for Parkinson's disease." \textit{Movement Disorders}, 2015;30(12):1591--1601.
\bibitem{muscle_rigidity:ref3} Olanow CW, Stern MB, Sethi K. "The scientific and clinical basis for the treatment of Parkinson disease." \textit{Neurology}, 2009;72(21 Suppl 4):S1--S136.
\bibitem{muscle_rigidity:ref4} Fahn S. "Classification of movement disorders." \textit{Movement Disorders}, 2011;26(6):947--957.
\bibitem{muscle_rigidity:ref5} Samii A, Nutt JG, Ransom BR. "Parkinson's disease." \textit{Lancet}, 2004;363(9423):1783--1793.
\bibitem{muscle_rigidity:ref6} Poewe W, Seppi K, Tanner CM, et al. "Parkinson disease." \textit{Nature Reviews Disease Primers}, 2017;3:17013.

\end{thebibliography}
